% Ch7 WS3 -- periodic trends, successive IE, PES, plus an FRQ. Block 4.
\documentclass{apchem-worksheet}

\apcunitword{Chapter}
\apcunit{7}
\apcunittitle{Atomic Structure and Periodicity}
\apcdoctitle{Worksheet 3 \textbullet\ Trends, Ionization Energy \& PES}
\apcsubtitle{Zumdahl \S7.12 \textbullet\ Coulombic reasons only --- no octet language}

\begin{document}
\wsheader[Every trend explanation must use nuclear charge, distance, or
shielding. ``It wants a full shell'' and ``it is more stable'' earn zero on
the AP exam, so they earn zero here.]

\problem[]{Rank each set and give the Coulombic reason:}
\begin{parts}
  \item Atomic radius, largest first: Li, Na, K
        \answerlines[2]{K $>$ Na $>$ Li --- each successive element adds a
        shell, placing valence electrons farther from the nucleus.}
  \item Atomic radius, largest first: Na, Mg, Al
        \answerlines[2]{Na $>$ Mg $>$ Al --- same shell, but rising nuclear
        charge increases $Z_{\text{eff}}$ and pulls the electron cloud
        inward.}
  \item First ionization energy, largest first: F, O, N
        \answerlines[2]{F $>$ N $>$ O --- generally increases across the
        period, but O dips below N because O's fourth 2$p$ electron is
        paired and the added repulsion makes it easier to remove.}
\end{parts}

\problem[]{Explain each period-2 anomaly:}
\begin{parts}
  \item Why is $I_1$ of boron less than that of beryllium?
        \answerlines[2]{Boron's outermost electron occupies a 2$p$ orbital,
        which is higher in energy and less penetrating than beryllium's
        filled 2$s$.}
  \item Why is $I_1$ of oxygen less than that of nitrogen?
        \answerlines[2]{Nitrogen's 2$p$ electrons are all unpaired; oxygen's
        fourth 2$p$ electron must pair up, and the electron--electron
        repulsion makes it easier to remove.}
\end{parts}

\problem[]{Successive ionization energies (kJ/mol). Identify the group and
justify from the data:}
\begin{parts}
  \item 496, 4562, 6910, 9543 $\to$
        \answerblank[5cm]{group 1 --- jump after $I_1$}
  \item 786, 1577, 3232, 4356, 16{,}091 $\to$
        \answerblank[5cm]{group 14 --- jump after $I_4$}
  \item 578, 1817, 2745, 11{,}577 $\to$
        \answerblank[6.2cm]{group 13 --- jump after $I_3$ (Al)}
\end{parts}

\problem[]{Explain why $I_2$ is always greater than $I_1$, regardless of
element. \answerlines[3]{The second electron is removed from a positively
charged ion rather than a neutral atom. With one fewer electron and the same
nuclear charge, the remaining electrons experience less mutual repulsion and
a stronger net attraction --- Coulomb's law --- so more energy is required.}}

\problem[]{Photoelectron spectroscopy. A spectrum shows peaks at 273, 26.8,
20.2, 2.44, and 1.25 MJ/mol with relative areas $2:2:6:2:5$.}
\begin{parts}
  \item Write the electron configuration.
        \answerblank[5.5cm]{$1s^2 2s^2 2p^6 3s^2 3p^5$}
  \item Identify the element, citing your evidence.
        \answerblank[6.0cm]{Cl --- areas total 17 electrons}
  \item Which peak corresponds to the electrons removed first on
        ionization? \answerblank[4.5cm]{1.25 MJ/mol (3$p$)}
  \item Why does the 2$p$ peak (20.2) lie at lower binding energy than the
        2$s$ peak (26.8), given both are in shell 2?
        \answerlines[2]{2$s$ penetrates closer to the nucleus, so it is
        shielded less and held more tightly; 2$p$ is shielded slightly more
        and therefore bound less strongly.}
\end{parts}

\problem[]{\textbf{FRQ (10 points).} Answer all parts.}
\begin{parts}
  \item The complete PES spectrum of an element shows peaks at 208, 18.7,
        13.5, 1.95, and 1.06 MJ/mol with relative areas $2:2:6:2:3$.
        Determine the electron configuration and identify the element.
        \answerlines[2]{$1s^2 2s^2 2p^6 3s^2 3p^3$; areas total 15
        electrons $\to$ phosphorus.}
  \item Explain what the \emph{position} and the \emph{area} of a PES peak
        each represent. \answerlines[2]{Position gives the binding energy of
        the electrons in that subshell; area is proportional to the number
        of electrons occupying it.}
  \item Predict two specific ways sulfur's spectrum would differ from this
        one. \answerlines[2]{Every peak shifts to slightly higher binding
        energy, because sulfur has one more proton at essentially the same
        shielding; and the lowest-energy peak's area increases from 3 to 4.}
  \item The 1$s$ peak sits at 208 MJ/mol while the 3$p$ peak sits at 1.06
        --- a factor of nearly 200. Explain this enormous gap.
        \answerlines[3]{1$s$ electrons are extremely close to the nucleus
        and essentially unshielded, so they feel nearly the full $+15$
        charge. The 3$p$ electrons are far away and shielded by all ten
        core electrons, so their effective nuclear charge is a small
        fraction of that.}
  \item Successive ionization energies for this element are 1012, 1907,
        2914, 4964, 6274, and 21{,}267 kJ/mol. Explain how these data
        independently confirm your identification in part (a).
        \answerlines[3]{The first large jump occurs between $I_5$ and
        $I_6$, showing five valence electrons before the core is reached.
        That places the element in group 15, consistent with phosphorus and
        with the $3s^2 3p^3$ configuration from the PES areas.}
\end{parts}

\begin{apcnote}[Rubric --- score yourself, 10 points]
\textbf{(a) 2 pts} 1 configuration matching all five areas, 1 phosphorus
tied to the 15-electron total \textbullet\ \textbf{(b) 2 pts} 1 position,
1 area \textbullet\ \textbf{(c) 2 pts} 1 shift to higher BE with the proton
reason, 1 area change \textbullet\ \textbf{(d) 2 pts} 1 says 1$s$ is
unshielded/close, 1 says 3$p$ is shielded and distant --- ``it's in a
different shell'' alone earns 1 \textbullet\ \textbf{(e) 2 pts} 1 locates
the jump after $I_5$, 1 links five valence electrons to group 15.
\end{apcnote}

\begin{spiralstrip}{Chapter 7 \textbullet\ blocks 1--3}
  \item Configuration of \ce{Fe^3+}: \answerblank[3.4cm]{$[\ce{Ar}]3d^5$}
  \item Energy of a \SI{600}{\nano\meter} photon:
        \answerblank[3.2cm]{\SI{3.31e-19}{\joule}}
  \item Unpaired electrons in ground-state O: \answerblank[1.6cm]{2}
  \item Which fills first, 4$s$ or 3$d$? \answerblank[1.8cm]{4$s$}
  \item Actual configuration of \ce{Cu}:
        \answerblank[3.6cm]{$[\ce{Ar}]4s^1 3d^{10}$}
\end{spiralstrip}

\end{document}
